% =====================================================================
% semana2.tex — Instalación de Oracle XE 21c y creación de BAHIA.inf
% Requiere los paquetes y macros de estilo.tex.
% Si tu documento es de clase article, cambia \chapter por \section y
% baja un nivel el resto (\section -> \subsection, etc.).
% =====================================================================

% Los cabezazos 1-3 corresponden a la semana 1 (ver el log del final).
\setcounter{cabezazo}{3}

\chapter{Instalación de Oracle XE 21c y creación de BAHIA.inf}
\label{cap:semana2}

En esta fase se instala Oracle Database 21c Express Edition sobre la máquina
\cmd{delfos.inf} (Ubuntu Server 26.04.1 LTS) preparada en la semana anterior
según la arquitectura OFA, y se crea la base de datos \textbf{BAHIA.inf} con los
parámetros que pide el enunciado. Al final se instalan los esquemas de ejemplo y
se configura el arranque automático.

\begin{table}[H]
\centering\small
\begin{tabular}{@{}ll@{}}
\toprule
\textbf{Parámetro} & \textbf{Valor} \\
\midrule
Software & Oracle Database 21c XE (21.3.0.0.0) \\
\texttt{ORACLE\_BASE} / \texttt{ORACLE\_HOME} & \cmd{/opt/oracle} / \cmd{/opt/oracle/product/21c/dbhomeXE} \\
Nombre global & \cmd{BAHIA.inf} (\cmd{db_name=BAHIA}, \cmd{db_domain=inf}) \\
SID & \cmd{XE} \\
Juego de caracteres & \cmd{AL32UTF8} (nacional \cmd{AL16UTF16}) \\
Tamaño de bloque & 8192 bytes (8 KB) \\
Arquitectura & CDB con una PDB: \cmd{BAHIAPDB} \\
Ficheros de datos & \cmd{/opt/oracle/oradata} \\
Memoria & SGA automática, 1536 MB \\
Esquemas de ejemplo & \cmd{HR} en \cmd{BAHIAPDB} \\
\bottomrule
\end{tabular}
\caption{Resumen de la configuración de BAHIA.inf}
\label{tab:bahia}
\end{table}

Como regla de trabajo durante toda la fase: \textbf{lo que necesita
\cmd{sudo} se hace con el usuario administrador \cmd{sercal}; lo propio de Oracle
(\cmd{sqlplus}, \cmd{lsnrctl}, \cmd{dbca}) con el usuario \cmd{oracle}}, que no
tiene privilegios de administración.

% ---------------------------------------------------------------------
\section{Descarga y verificación del software}
\label{sec:descarga}

El paquete se descarga de la web de Oracle (se necesita una cuenta gratuita).
Para Linux solo se distribuye como RPM para Oracle Linux 8; es la versión
\emph{XE} más reciente (la siguiente generación, 23ai, se distribuye como
\emph{Oracle Database Free} y no permite crear la base de datos con DBCA como
pide el enunciado).

\captura{oraclexedwnld.png}{Página de descarga de Oracle Database 21c XE}{fig:descarga}

Se comprueba la integridad del fichero comparando su SHA-256 con el publicado por
Oracle:

\begin{lstlisting}[style=bash]
sha256sum oracle-database-xe-21c-1.0-1.ol8.x86_64.rpm
# f8357b432de33478549a76557e8c5220ec243710ed86115c65b0c2bc00a848db
\end{lstlisting}

\captura{sha256oracle.png}{El hash calculado coincide con el publicado por Oracle}{fig:sha256}

El RPM se copia a la máquina virtual desde el anfitrión:

\begin{lstlisting}[style=bash]
scp oracle-database-xe-21c-1.0-1.ol8.x86_64.rpm sercal@delfos.inf:~/
\end{lstlisting}

\captura{scporacle.png}{Copia del paquete (2,2 GB) a \texttt{delfos.inf}}{fig:scp}

% ---------------------------------------------------------------------
\section{Conversión del RPM a paquete Debian}
\label{sec:alien}

Ubuntu usa paquetes \cmd{.deb}, así que el RPM se convierte con \cmd{alien}
conservando los scripts de instalación (\cmd{--scripts}):

\begin{lstlisting}[style=bash]
sudo alien --scripts -d oracle-database-xe-21c-1.0-1.ol8.x86_64.rpm
\end{lstlisting}

Los avisos \cmd{NOKEY} que aparecen repetidos solo indican que la clave GPG de
Oracle no está importada en el sistema; no son errores.

\captura{alien.png}{Primer intento de conversión con \texttt{alien}}{fig:alien}

La conversión resultó ser el paso más delicado por falta de espacio, primero en
el disco y después en \cmd{/tmp}.

\begin{cabezazo}{\texttt{alien} falla con «No space left on device»}{cab:disco}
\problema La conversión aborta copiando las librerías de
\cmd{dbhomeXE/lib} con \cmd{cp: error copying ... No space left on device} y
\cmd{make: *** [debian/rules:7: binary] Error 2}.

\causa \cmd{alien} descomprime el RPM y lo vuelve a copiar a
\cmd{debian/} antes de empaquetar, por lo que necesita unos 20 GB libres. El
instalador de Ubuntu Server solo asignó 19 GB al volumen lógico raíz
(\cmd{ubuntu-lv}) y dejó otros 19 GB libres en el grupo de volúmenes:

\begin{lstlisting}[style=fichero]
/dev/mapper/ubuntu--vg-ubuntu--lv   19G  9,6G  8,0G  55% /
  VG        #PV #LV #SN Attr   VSize   VFree
  ubuntu-vg   1   1   0 wz--n- <38,00g 19,00g
\end{lstlisting}

\solucion Primero se asigna al volumen raíz el espacio libre del grupo y,
para tener margen para la base de datos, se amplía además el disco virtual en
30 GB desde el anfitrión:

\begin{lstlisting}[style=bash]
# En la VM: usar el espacio libre del VG
sudo lvextend -r -l +100%FREE /dev/ubuntu-vg/ubuntu-lv

# En el anfitrión, con la VM apagada: ampliar el qcow2
sudo qemu-img resize /var/lib/libvirt/images/delfos.qcow2 +30G

# En la VM: crecer partición, volumen físico y volumen lógico
sudo apt install -y cloud-guest-utils
sudo growpart /dev/vda 3
sudo pvresize /dev/vda3
sudo lvextend -r -l +100%FREE /dev/ubuntu-vg/ubuntu-lv
df -h /
\end{lstlisting}

El sistema de ficheros raíz pasa a 67 GB con 55 GB libres
(figuras~\ref{fig:qemuimg} y~\ref{fig:lvextend}).
\end{cabezazo}

\captura[0.6\linewidth]{cabezazoalien.png}{Ampliación del disco virtual con \texttt{qemu-img resize}}{fig:qemuimg}

\captura{solucioncabezazoalien.png}{\texttt{growpart}, \texttt{pvresize} y \texttt{lvextend}: el volumen raíz pasa de 38 a 68 GiB}{fig:lvextend}

\begin{cabezazo}{\texttt{zstd write error} con 55 GB libres}{cab:tmpfs}
\problema Con espacio de sobra en el disco, \cmd{alien} vuelve a fallar al
final con \cmd{dpkg-deb (subproceso): comprimiendo miembro tar: zstd write
error: No space left on device}, y el \cmd{.deb} generado pesa solo 521 KB.

\causa \cmd{dpkg-deb} comprime el paquete en un fichero temporal en
\cmd{/tmp}, y en Ubuntu 26.04 \cmd{/tmp} es un \emph{tmpfs} en memoria RAM de
solo 1,7 GB (la mitad de la RAM de la VM):

\begin{lstlisting}[style=fichero]
$ df -h /tmp
Filesystem      Size  Used Avail Use% Mounted on
tmpfs           1,7G     0  1,7G   0% /tmp
\end{lstlisting}

\solucion Indicar un directorio temporal en disco con la variable
\cmd{TMPDIR}. Se usa \cmd{sudo env} porque \cmd{sudo} descarta las variables de
entorno del usuario:

\begin{lstlisting}[style=bash]
mkdir -p ~/tmp-alien
sudo env TMPDIR=/home/sercal/tmp-alien alien --scripts -d \
     oracle-database-xe-21c-1.0-1.ol8.x86_64.rpm
ls -lh *.deb
\end{lstlisting}

Los numerosos mensajes \cmd{dpkg-shlibdeps: error: cannot find library ...}
que aparecen durante la conversión no son el problema: \cmd{alien} los ignora
(\cmd{Error 25 (ignored)}) porque son librerías que Oracle incluye en su propio
\texttt{ORACLE\_HOME}.
\end{cabezazo}

\captura{debcreado.png}{Paquete \texttt{oracle-database-xe-21c\_1.0-2\_amd64.deb} generado correctamente (2,2 GB)}{fig:deb}

% ---------------------------------------------------------------------
\section{Instalación del paquete}
\label{sec:dpkg}

Antes de instalar se crea el enlace que Oracle espera para \cmd{libaio}
(en Ubuntu 26.04 la librería se llama \cmd{libaio.so.1t64}) y se comprueba que
el \texttt{ORACLE\_BASE} pertenece al usuario \cmd{oracle}:

\begin{lstlisting}[style=bash]
sudo ln -sf /usr/lib/x86_64-linux-gnu/libaio.so.1t64 \
            /usr/lib/x86_64-linux-gnu/libaio.so.1
sudo chown oracle:oinstall /opt/oracle /opt/oraInventory
sudo dpkg -i oracle-database-xe-21c_1.0-2_amd64.deb
\end{lstlisting}

\begin{cabezazo}{\texttt{awk: command not found} en el script de preinstalación}{cab:awk}
\problema \cmd{dpkg -i} se detiene en el \emph{preinst}:
\cmd{/var/lib/dpkg/tmp.ci/preinst: line 188: awk: command not found}.

\causa Siguiendo una receta pensada para distribuciones antiguas se ejecutó
\cmd{sudo ln -sf /usr/bin/awk /bin/awk}. En Ubuntu moderno \cmd{/bin} es un
enlace a \cmd{/usr/bin} (\emph{usrmerge}), así que el comando sobrescribió
\cmd{/usr/bin/awk} con un enlace que apunta a sí mismo
(\cmd{/usr/bin/awk -> /usr/bin/awk}) y \cmd{awk} dejó de existir en todo el
sistema.

\solucion Borrar el enlace roto y reinstalar \cmd{mawk}, que regenera el
grupo de alternativas de \cmd{awk}:

\begin{lstlisting}[style=bash]
sudo rm /usr/bin/awk
sudo apt install --reinstall mawk
ls -l /usr/bin/awk      # /usr/bin/awk -> /etc/alternatives/awk
awk --version
\end{lstlisting}

El aviso \cmd{ln: failed to create symbolic link '/bin/awk': Already exists}
que muestra después el instalador de Oracle es inofensivo.
\textbf{No hay que crear ningún enlace para \cmd{awk} en Ubuntu.}
\end{cabezazo}

\captura{cabezazoawk.png}{El enlace circular de \texttt{awk} y su reparación reinstalando \texttt{mawk}}{fig:awk}

\begin{cabezazo}{\texttt{ORACLE\_BASE} no pertenece al usuario \texttt{oracle}}{cab:oraclebase}
\problema En la misma ejecución, el preinstalador muestra
\cmd{[SEVERE] The install cannot proceed because ORACLE_BASE directory
(/opt/oracle) is not owned by "oracle" user}.

\causa El paquete comprueba la propiedad de \cmd{/opt/oracle} antes de
desempaquetar (y lo hace usando \cmd{awk}, por lo que el problema anterior
agravaba este).

\solucion \cmd{sudo chown -R oracle:oinstall /opt/oracle /opt/oraInventory}
\emph{antes} de instalar. Una vez instalado Oracle \textbf{no debe repetirse el
\cmd{chown -R}}: ver el cabezazo~\ref{cab:oradism}.
\end{cabezazo}

Con \cmd{awk} reparado, la instalación termina correctamente:

\captura{solawk.png}{\texttt{Oracle home installed successfully and ready to be configured}}{fig:dpkgok}

El instalador sugiere ejecutar \cmd{/etc/init.d/oracle-xe-21c configure}, pero
\textbf{no se ejecuta}: ese script crearía una base de datos \cmd{XE} con valores
por defecto, mientras que BAHIA debe crearse con DBCA y los parámetros del
enunciado.

Se toma una instantánea de la VM en este punto y se borran los paquetes
descargados (los originales se conservan en el directorio \cmd{software/} del
anfitrión):

\begin{lstlisting}[style=bash]
# Anfitrión, con la VM apagada
virsh -c qemu:///system snapshot-create-as delfos --name "s2-xe-instalado" \
      --description "Oracle XE 21c instalado, sin BAHIA"
# VM
rm -f ~/oracle-database-xe-21c-1.0-1.ol8.x86_64.rpm
sudo rm -f ~/oracle-database-xe-21c_1.0-2_amd64.deb
\end{lstlisting}

\captura{snapshot2.png}{Instantánea \texttt{s2-xe-instalado}}{fig:snap2}
\captura[0.7\linewidth]{limpiezapostsnpcht.png}{Limpieza de paquetes: 51 GB libres}{fig:limpieza}

% ---------------------------------------------------------------------
\section{Entorno del usuario \texttt{oracle} y listener}
\label{sec:entorno}

Las variables de entorno definidas en la semana 1 en
\cmd{/home/oracle/.bash_profile} apuntan ya al \texttt{ORACLE\_HOME} instalado.
Se comprueba además que los binarios no tienen librerías sin resolver:

\begin{lstlisting}[style=bash]
sudo su - oracle
echo $ORACLE_HOME $ORACLE_SID
which sqlplus dbca lsnrctl
ldd $ORACLE_HOME/bin/oracle  | grep "not found"
ldd $ORACLE_HOME/bin/sqlplus | grep "not found"
sqlplus -v
\end{lstlisting}

\captura[0.6\linewidth]{entornousuariooracle.png}{Entorno de \texttt{oracle} correcto y sin librerías ausentes}{fig:entorno}

Al no haber usado el \cmd{configure} de XE, el listener se crea a mano. Escucha
por TCP en \cmd{delfos.inf:1521} (el nombre se resuelve con \cmd{/etc/hosts}) y
por IPC:

\begin{lstlisting}[style=fichero,caption={\texttt{\$ORACLE\_HOME/network/admin/listener.ora}}]
LISTENER =
  (DESCRIPTION_LIST =
    (DESCRIPTION =
      (ADDRESS = (PROTOCOL = TCP)(HOST = delfos.inf)(PORT = 1521))
      (ADDRESS = (PROTOCOL = IPC)(KEY = EXTPROC1521))
    )
  )
\end{lstlisting}

\begin{lstlisting}[style=bash]
lsnrctl start
lsnrctl status
\end{lstlisting}

\captura{listener.png}{Listener arrancado. Aún «no soporta ningún servicio» porque no existe ninguna base de datos}{fig:listener}

% ---------------------------------------------------------------------
\section{Creación de BAHIA.inf con DBCA}
\label{sec:dbca}

\subsection{Intento con la interfaz gráfica}

La máquina es un Ubuntu Server sin entorno gráfico, así que la idea era lanzar
DBCA reenviando su ventana al escritorio del anfitrión por SSH (X11
\emph{forwarding}). Se instalaron las dependencias gráficas mínimas en la VM:

\begin{lstlisting}[style=bash]
sudo apt install -y xauth libxrender1 libxtst6 libxi6 fontconfig
\end{lstlisting}

\captura{x11.png}{Dependencias para el reenvío X11}{fig:x11deps}

\begin{cabezazo}{DBCA: «DISPLAY not set»}{cab:display}
\problema Al ejecutar \cmd{dbca} tras \cmd{sudo su - oracle} aparece
\cmd{DISPLAY not set. Set DISPLAY environment variable, then re-run.}

\causa La variable \cmd{DISPLAY} solo la crea una conexión \cmd{ssh -X} o
\cmd{ssh -Y} hecha desde el anfitrión gráfico; una sesión abierta con
\cmd{su} desde la consola de la VM no tiene pantalla a la que enviar la
ventana. Además, al pegar de golpe un bloque que empieza por
\cmd{sudo su - oracle}, las líneas siguientes se pierden mientras \cmd{sudo}
pide la contraseña.

\solucion Conectarse con \cmd{ssh -Y oracle@delfos.inf} desde el escritorio
del anfitrión, o copiar la \emph{cookie} X11 de \cmd{sercal} a \cmd{oracle}
(\cmd{xauth list $DISPLAY} / \cmd{xauth add ...}) ejecutando las órdenes de una
en una.
\end{cabezazo}

\begin{cabezazo}{DBCA: «Invalid MIT-MAGIC-COOKIE-1 key»}{cab:cookie}
\problema Con \cmd{DISPLAY=localhost:11.0} y la \emph{cookie} correctamente
copiada (aparece en \cmd{xauth list}), DBCA falla con
\cmd{Invalid MIT-MAGIC-COOKIE-1 key} y
\cmd{java.lang.NoClassDefFoundError: Could not initialize class
sun.awt.X11.XToolkit}.

\causa El rechazo viene del lado del anfitrión: al conectar, SSH avisa
\cmd{Warning: No xauth data; using fake authentication data for X11
forwarding}. El escritorio del anfitrión (Wayland/XWayland) no ofrece una
\emph{cookie} X11 válida, así que la que llega a la VM es falsa y el servidor
gráfico rechaza la ventana, tanto con \cmd{-X} como con \cmd{-Y}.

\solucion Se abandona la interfaz gráfica y se usa DBCA en \textbf{modo
silencioso} (\cmd{-silent}), que ofrece las mismas opciones, incluida la
generación de los scripts de creación, y queda mejor documentado en el
repositorio (\cmd{scripts/03-dbca-bahia.sh}).
\end{cabezazo}

\captura[0.75\linewidth]{cabezazox11.png}{«No xauth data; using fake authentication data» al conectar con \texttt{ssh -Y}}{fig:x11a}
\captura[0.75\linewidth]{cabezazox11_2.png}{La \emph{cookie} de \texttt{sercal} y el bloque pegado que se pierde tras \texttt{sudo su}}{fig:x11b}
\captura{rendicionx11.png}{Con la \emph{cookie} copiada, DBCA sigue fallando: se pasa al modo silencioso}{fig:x11c}

\subsection{Generación de los scripts de creación}

El enunciado pide activar la opción de DBCA que genera los ficheros de comandos
de creación. En modo silencioso es la operación \cmd{-generateScripts}. Primero
se comprueba el nombre de la plantilla de XE:

\begin{lstlisting}[style=bash]
ls $ORACLE_HOME/assistants/dbca/templates/
# XE_Database.dbc  XE_Seed_Database.ctl  XE_Seed_Database.dfb  pdbseed.dfb  pdbseed.xml
mkdir -p /opt/oracle/admin/BAHIA/scripts
\end{lstlisting}

\begin{lstlisting}[style=bash]
dbca -silent -generateScripts \
  -templateName XE_Database.dbc \
  -gdbName BAHIA.inf -sid XE \
  -characterSet AL32UTF8 -nationalCharacterSet AL16UTF16 \
  -createAsContainerDatabase true -numberOfPDBs 1 -pdbName BAHIAPDB \
  -datafileDestination /opt/oracle/oradata -storageType FS \
  -memoryMgmtType AUTO_SGA -totalMemory 1536 \
  -emConfiguration NONE \
  -scriptDest /opt/oracle/admin/BAHIA/scripts
\end{lstlisting}

El tamaño de bloque no se indica porque la plantilla de XE lo fija en 8192 bytes,
justo el valor que pide el enunciado.

\begin{cabezazo}{\texttt{[DBT-12513]} Esquemas de ejemplo no soportados}{cab:dbt12513}
\problema Con la opción \cmd{-sampleSchema true}, DBCA responde
\texttt{[FATAL] [DBT-12513] El esquema de muestra no está soportado en la versión
del directorio raíz de Oracle 21.0.0.0.0.}

\causa Desde Oracle 21c, DBCA ya no incluye los esquemas de ejemplo; Oracle
los publica aparte en GitHub (\cmd{oracle-samples/db-sample-schemas}).

\solucion Crear la base de datos sin \cmd{-sampleSchema} e instalar los
esquemas después desde el repositorio oficial (sección~\ref{sec:hr}).
\end{cabezazo}

\begin{cabezazo}{\texttt{-generateScripts} termina con \texttt{NullPointerException}}{cab:npe}
\problema Tras mostrar todos los porcentajes hasta el 100\,\%, el comando
acaba con \cmd{[FATAL] java.lang.NullPointerException}. Parecía que la base de
datos podía haberse creado a medias.

\causa Es un fallo de DBCA al terminar la generación de scripts en modo
silencioso. El progreso es simulado: en el log todos los pasos ocurren en el
mismo segundo (20:47:51), no hay proceso \cmd{pmon}, \cmd{/opt/oracle/oradata}
está vacío y \cmd{/etc/oratab} no tiene entradas.

\solucion Ninguna: los 17 scripts se generan completos en
\cmd{/opt/oracle/admin/BAHIA/scripts}. Se verifican y se continúa con
\cmd{-createDatabase}.
\end{cabezazo}

\captura{scripts1.png}{\texttt{DBT-12513} con \texttt{-sampleSchema} y la generación de scripts sin él}{fig:gen}
\captura[0.8\linewidth]{cabezazoscript1.png}{Comprobación: no hay instancia ni ficheros de datos, pero sí los scripts generados}{fig:genchk}

Los scripts generados confirman los parámetros del enunciado
(\cmdb{db_name="BAHIA"}, \cmdb{db_domain="inf"}, \cmdb{db_block_size=8192},
\cmdb{AL32UTF8}) y no contienen contraseñas en claro: las piden en tiempo de
ejecución mediante las variables de sustitución \cmdb{&&sysPassword} y
\cmdb{&&systemPassword}. Por eso pueden subirse tal cual al directorio
\cmdb{db-creation-scripts/} del repositorio.

\begin{lstlisting}[style=bash]
cd /opt/oracle/admin/BAHIA/scripts
grep -E "db_name|db_domain|db_block_size" init.ora
grep -i -E "characterset|AL32UTF8" *.sql
grep -i -n "password\|identified by" *.sql *.sh
\end{lstlisting}

\captura{cabezazoscript2.png}{Parámetros de BAHIA en \texttt{init.ora} y contraseñas sustituidas por variables}{fig:scriptsgrep}

\begin{table}[H]
\centering\small
\begin{tabular}{@{}>{\raggedright\arraybackslash}p{0.36\linewidth}p{0.58\linewidth}@{}}
\toprule
\textbf{Fichero} & \textbf{Función} \\
\midrule
\cmdb{XE.sh} & Script principal: crea directorios y lanza \cmdb{XE.sql}. \\
\cmdb{XE.sql} & Pide las contraseñas y encadena el resto de scripts. \\
\cmdb{init.ora} & Parámetros de inicialización de la instancia. \\
\cmdb{cloneDBCreation.sql} & Crea la CDB a partir de la plantilla \emph{seed}. \\
\cmdb{CloneRmanRestore.sql}, \cmdb{rman*.sql} & Restauran con RMAN los ficheros de la plantilla. \\
\cmdb{plug_PDBSeed.sql}, \cmdb{plug_BAHIAPDB.sql} & Conectan \cmdb{PDB$SEED} y crean \cmdb{BAHIAPDB}. \\
\cmdb{postScripts.sql}, \cmdb{postDBCreation.sql} & Compilan objetos y finalizan la configuración. \\
\cmdb{lockAccount.sql} & Bloquea las cuentas internas no usadas. \\
\bottomrule
\end{tabular}
\caption{Principales scripts de creación generados por DBCA}
\label{tab:scripts}
\end{table}

\subsection{Creación de la base de datos}

La creación se lanza dentro de \cmd{tmux} para que no se interrumpa si se corta
la conexión SSH. DBCA pide por teclado las contraseñas de \cmd{SYS},
\cmd{SYSTEM} y \cmd{PDBADMIN}, de modo que no quedan en el historial:

\begin{lstlisting}[style=bash]
dbca -silent -createDatabase \
  -templateName XE_Database.dbc \
  -gdbName BAHIA.inf -sid XE \
  -characterSet AL32UTF8 -nationalCharacterSet AL16UTF16 \
  -createAsContainerDatabase true -numberOfPDBs 1 -pdbName BAHIAPDB \
  -datafileDestination /opt/oracle/oradata -storageType FS \
  -memoryMgmtType AUTO_SGA -totalMemory 1536 \
  -emConfiguration NONE
\end{lstlisting}

\begin{cabezazo}{\texttt{[DBT-10011]} Permisos incorrectos de \texttt{oradism}}{cab:oradism}
\problema La creación se detiene con
\cmd{[FATAL] [DBT-10011] Se han detectado una propiedad o permisos incorrectos
para el archivo (/opt/oracle/product/21c/dbhomeXE/bin/oradism)}, acompañado de
\cmd{PRVG-11960} (falta el bit SUID) y \cmd{PRVG-2031} (se esperaba
\cmd{root(0)} y se encontró \cmd{oracle(54321)}).

\causa \cmd{oradism} gestiona la memoria compartida de la instancia y debe
pertenecer a \cmd{root} con el bit SUID. Los \cmd{chown -R oracle:oinstall
/opt/oracle} ejecutados tras la instalación le cambiaron el propietario y le
quitaron ese bit.

\solucion Ejecutar como \cmd{root} el script \cmd{root.sh} del
\texttt{ORACLE\_HOME}, que restablece los permisos especiales:

\begin{lstlisting}[style=bash]
sudo /opt/oracle/product/21c/dbhomeXE/root.sh
ls -l /opt/oracle/product/21c/dbhomeXE/bin/oradism
# -rwsr-x--- 1 root oinstall 1867560 Aug 18  2021 .../oradism
\end{lstlisting}

\cmd{root.sh} deja su salida en un log que \cmd{sercal} solo puede leer con
\cmd{sudo} (de ahí un \cmd{Permission denied} al hacer \cmd{cat} sin él).
\textbf{A partir de aquí no se vuelve a usar \cmd{chown -R} sobre
\cmd{/opt/oracle}.}
\end{cabezazo}

\begin{cabezazo}{\texttt{[DBT-06208]} Contraseñas débiles}{cab:pwd}
\problema DBCA muestra un aviso \cmd{DBT-06208} para \cmd{SYS},
\cmd{SYSTEM} y \cmd{PDBADMIN}.

\causa Oracle recomienda al menos 8 caracteres con mayúsculas, minúsculas y
dígitos, y desaconseja palabras como «oracle». Es un aviso, no un error.

\solucion Usar contraseñas robustas, ya que en la fase siguiente se abrirá
el puerto 1521 a la red para la conexión remota como administrador. Las
contraseñas se guardan en un gestor de contraseñas, nunca en el repositorio.
\end{cabezazo}

\captura{cabezazoscript3.png}{\texttt{DBT-06208} (aviso) y \texttt{DBT-10011} (error) en el primer intento}{fig:oradism}
\captura{script4.png}{Creación de la base de datos completada: \texttt{BAHIA.inf}, SID \texttt{XE}}{fig:dbcaok}

Esta vez cada fase tarda minutos, y la creación termina con
\cmd{Nombre de la Base de Datos Global: BAHIA.inf} e
\cmd{Identificador del Sistema (SID): XE}.

\subsection{Verificación}

\begin{lstlisting}[style=sql]
-- sqlplus / as sysdba
SELECT * FROM global_name;
SELECT value FROM nls_database_parameters WHERE parameter='NLS_CHARACTERSET';
SHOW PARAMETER db_block_size
SELECT instance_name, status FROM v$instance;
SHOW PDBS
\end{lstlisting}

\captura{porfinterminoscripts.png}{BAHIA.INF, AL32UTF8, bloque de 8192, instancia XE abierta y BAHIAPDB en lectura/escritura}{fig:verif}

La base de datos se registra automáticamente en el listener, que ahora ofrece
los servicios \cmd{BAHIA.inf} (CDB) y \cmd{bahiapdb.inf} (PDB). DBCA también
añade la entrada de la instancia en \cmd{/etc/oratab}, todavía con arranque
automático desactivado (\cmd{N}).

\captura[0.8\linewidth]{listenercheck.png}{Servicios registrados en el listener y entrada en \texttt{/etc/oratab}}{fig:lsnrchk}

Para hacer una instantánea coherente se detienen la base de datos y el listener,
se apaga la VM y se crea la instantánea \cmd{s2-bahia-creada}:

\captura[0.8\linewidth]{presnapshot3.png}{Parada ordenada con \texttt{SHUTDOWN IMMEDIATE}}{fig:shutdown}
\captura{snapshot3.png}{Instantánea \texttt{s2-bahia-creada}}{fig:snap3}

Tras volver a arrancar la VM, la base de datos se levanta a mano:

\captura[0.75\linewidth]{startup.png}{\texttt{lsnrctl start} y \texttt{STARTUP}: SGA de 1,5 GB, base de datos abierta}{fig:startup}

\begin{cabezazo}[unbreakable]{La PDB queda en \texttt{MOUNTED} tras reiniciar}{cab:pdb}
\problema Después del \cmd{STARTUP}, \cmd{SHOW PDBS} muestra
\cmd{BAHIAPDB} en estado \cmd{MOUNTED}, no accesible.

\causa Por defecto las PDB no se abren al arrancar la CDB.

\solucion Abrirla y guardar su estado para que se abra sola en cada
arranque:

\begin{lstlisting}[style=sql]
ALTER PLUGGABLE DATABASE BAHIAPDB OPEN;
ALTER PLUGGABLE DATABASE BAHIAPDB SAVE STATE;
\end{lstlisting}
\end{cabezazo}

\captura[0.75\linewidth]{showpdbs.png}{\texttt{BAHIAPDB} abierta en \texttt{READ WRITE} y con el estado guardado}{fig:pdb}

% ---------------------------------------------------------------------
\section{Esquemas de ejemplo}
\label{sec:hr}

Como DBCA 21c ya no los incluye (cabezazo~\ref{cab:dbt12513}), se descargan los
esquemas oficiales de Oracle y se instala el esquema \textbf{HR (Human
Resources)}, que es el que se usará para la migración a PostgreSQL:

\begin{lstlisting}[style=bash]
cd ~
wget https://github.com/oracle-samples/db-sample-schemas/archive/refs/heads/main.tar.gz \
     -O sample-schemas.tar.gz
tar xzf sample-schemas.tar.gz
cd db-sample-schemas-main/human_resources
sqlplus system@//delfos.inf:1521/bahiapdb.inf
\end{lstlisting}

\begin{lstlisting}[style=sql]
SHOW CON_NAME        -- debe mostrar BAHIAPDB
@hr_install.sql
\end{lstlisting}

El instalador pide la contraseña de \cmd{HR}, el tablespace (\cmd{USERS}) y si
debe sobrescribir el esquema (\cmd{YES}).

\begin{cabezazo}{\texttt{ORA-65096} al crear el usuario HR}{cab:ora65096}
\problema El primer intento de \cmd{hr_install.sql} falla en
\cmd{CREATE USER hr ...} con \texttt{ORA-65096: nombre de usuario o rol común no
válido}.

\causa El script se ejecutó conectado a la raíz del contenedor
(\cmd{CDB$ROOT}), donde solo pueden crearse usuarios comunes con prefijo
\texttt{C\#\#}. Los esquemas de ejemplo deben ir en la PDB.

\solucion Conectarse al servicio de la PDB (\cmd{bahiapdb.inf}, obtenido con
\cmd{lsnrctl status | grep -i bahia}) y comprobar con \cmd{SHOW CON_NAME} que se
está en \cmd{BAHIAPDB} antes de lanzar el script. Alternativa:
\cmd{ALTER SESSION SET CONTAINER = BAHIAPDB;} como \cmd{SYS}.
\end{cabezazo}

\begin{cabezazo}{\texttt{ORA-01017} al conectar como \texttt{SYSTEM}}{cab:ora01017}
\problema La primera conexión \cmd{sqlplus system@//delfos.inf:1521/bahiapdb.inf}
devuelve \texttt{ORA-01017: nombre de usuario/contraseña no válidos}.

\causa Contraseña mal tecleada (SQL*Plus no la muestra al escribirla).

\solucion SQL*Plus vuelve a pedir usuario y contraseña; al introducirlos
bien, la conexión funciona.
\end{cabezazo}

\captura{schema1.png}{Descarga de los esquemas oficiales y conexión a la PDB}{fig:hr1}
\captura[0.8\linewidth]{schema2.png}{Instalación de HR en \texttt{BAHIAPDB}}{fig:hr2}
\captura[0.8\linewidth]{schema3.png}{Verificación del instalador: todas las tablas con las filas esperadas}{fig:hr3}

\begin{table}[H]
\centering\small
\begin{tabular}{@{}lrr@{}}
\toprule
\textbf{Tabla} & \textbf{Esperadas} & \textbf{Cargadas} \\
\midrule
\cmd{regions} & 5 & 5 \\
\cmd{countries} & 25 & 25 \\
\cmd{departments} & 27 & 27 \\
\cmd{locations} & 23 & 23 \\
\cmd{employees} & 107 & 107 \\
\cmd{jobs} & 19 & 19 \\
\cmd{job_history} & 10 & 10 \\
\bottomrule
\end{tabular}
\caption{Verificación del esquema HR}
\label{tab:hr}
\end{table}

Por último, se comprueba el acceso con el propio usuario \cmd{HR}:

\captura[0.8\linewidth]{schemacheck.png}{Consulta a \texttt{employees} conectado como \texttt{HR}}{fig:hrchk}

% ---------------------------------------------------------------------
\section{Arranque automático}
\label{sec:arranque}

El paquete RPM de Oracle no crea en Ubuntu un servicio de arranque, así que se
define uno para \cmd{systemd} basado en las utilidades \cmd{dbstart} y
\cmd{dbshut}, que arrancan y detienen las instancias marcadas con \cmd{Y} en
\cmd{/etc/oratab} (y también el listener, al recibir el \texttt{ORACLE\_HOME}).

\begin{lstlisting}[style=bash]
sudo sed -i 's|^XE:/opt/oracle/product/21c/dbhomeXE:N|XE:/opt/oracle/product/21c/dbhomeXE:Y|' /etc/oratab
grep XE /etc/oratab
\end{lstlisting}

\captura[0.8\linewidth]{arranqueauto.png}{Activación del arranque automático en \texttt{/etc/oratab}}{fig:oratab}

\begin{lstlisting}[style=fichero,caption={\texttt{/etc/systemd/system/oracle-xe.service}}]
[Unit]
Description=Oracle Database XE 21c (BAHIA.inf) y listener
After=network-online.target
Wants=network-online.target

[Service]
Type=oneshot
RemainAfterExit=yes
User=oracle
Group=oinstall
Environment=ORACLE_HOME=/opt/oracle/product/21c/dbhomeXE
ExecStart=/opt/oracle/product/21c/dbhomeXE/bin/dbstart /opt/oracle/product/21c/dbhomeXE
ExecStop=/opt/oracle/product/21c/dbhomeXE/bin/dbshut /opt/oracle/product/21c/dbhomeXE
TimeoutSec=600

[Install]
WantedBy=multi-user.target
\end{lstlisting}

\captura[0.8\linewidth]{servicioracle.png}{Fichero de unidad del servicio}{fig:unit}

\begin{cabezazo}{«I'm sorry oracle. I'm afraid I can't do that»}{cab:sudo}
\problema Al ejecutar \cmd{sudo systemctl ...} desde la sesión del usuario
\cmd{oracle}, \cmd{sudo} lo rechaza y el servicio queda \cmd{disabled} e
\cmd{inactive (dead)}.

\causa El usuario \cmd{oracle} no pertenece al grupo \cmd{sudo}, y es lo
correcto: el propietario del software de Oracle no debe tener privilegios de
administración.

\solucion Cambiar al usuario administrador (\cmd{su - sercal}) y activar el
servicio desde ahí.
\end{cabezazo}

Antes de activarlo se detienen la base de datos y el listener arrancados a mano:

\captura[0.7\linewidth]{servicio2.png}{Parada manual previa a la activación del servicio}{fig:stopmanual}

\begin{lstlisting}[style=bash]
sudo systemctl daemon-reload
sudo systemctl enable --now oracle-xe
systemctl status oracle-xe
\end{lstlisting}

\captura{service3.png}{Servicio \texttt{enabled} y \texttt{active (exited)}: listener y procesos de la instancia \texttt{XE}}{fig:service}

La prueba definitiva es reiniciar la máquina y comprobar que todo vuelve solo,
incluida la PDB gracias al \cmd{SAVE STATE}:

\begin{lstlisting}[style=bash]
sudo reboot
# tras volver a entrar como sercal
systemctl status oracle-xe
ps -ef | grep [p]mon
sudo su - oracle -c 'echo "SHOW PDBS" | sqlplus -s / as sysdba'
\end{lstlisting}

\captura{servicecheck1.png}{Tras el reinicio, el servicio ha arrancado solo}{fig:reboot1}
\captura[0.85\linewidth]{servicecheck2.png}{Proceso \texttt{xe\_pmon\_XE} activo y \texttt{BAHIAPDB} en \texttt{READ WRITE} sin intervención manual}{fig:reboot2}

Con todo verificado se toma la instantánea final de la fase:

\captura{snapshot4.png}{Instantánea \texttt{s2-bahia-completa}}{fig:snap4}

% ---------------------------------------------------------------------
\section{Resumen de la fase}

\begin{table}[H]
\centering\small
\begin{tabular}{@{}p{0.55\linewidth}l@{}}
\toprule
\textbf{Requisito} & \textbf{Evidencia} \\
\midrule
Oracle XE 21c instalado según OFA & Fig.~\ref{fig:dpkgok}, \ref{fig:entorno} \\
Base de datos \cmd{BAHIA.inf}, SID \cmd{XE} & Fig.~\ref{fig:dbcaok}, \ref{fig:verif} \\
Juego de caracteres \cmd{AL32UTF8} & Fig.~\ref{fig:verif}, \ref{fig:scriptsgrep} \\
Bloque de 8 KB & Fig.~\ref{fig:verif}, \ref{fig:scriptsgrep} \\
Esquemas de ejemplo instalados & Fig.~\ref{fig:hr3}, tabla~\ref{tab:hr} \\
Scripts de creación generados & Fig.~\ref{fig:genchk}, tabla~\ref{tab:scripts} \\
Arranque automático & Fig.~\ref{fig:service}, \ref{fig:reboot2} \\
\bottomrule
\end{tabular}
\caption{Cumplimiento del enunciado en esta fase}
\end{table}

\begin{table}[H]
\centering\small
\begin{tabular}{@{}ll@{}}
\toprule
\textbf{Instantánea} & \textbf{Estado} \\
\midrule
\cmd{s1-preparada-pre-oracle} & Ubuntu + preparación OFA, sin Oracle \\
\cmd{s2-xe-instalado} & Oracle XE 21c instalado, sin base de datos \\
\cmd{s2-bahia-creada} & BAHIA.inf creada, sin esquemas de ejemplo \\
\cmd{s2-bahia-completa} & BAHIA.inf + HR + arranque automático \\
\bottomrule
\end{tabular}
\caption{Instantáneas de la máquina virtual}
\end{table}
